\begin{abstract}
Mathematical reasoning remains a formidable obstacle for language models, largely due to its intricacies and highly organized structure. In this study, we present DeepSeekMath~7B, an extension of DeepSeek-Coder-Base-v1.5 7B that undergoes continued pre-training with a dataset comprising 120B math-focused tokens extracted from Common Crawl, complemented by natural language and code content. DeepSeekMath~7B attains an outstanding 51.7\% accuracy on the MATH benchmark at the competition level, without the assistance of external toolkits or ensembling methods, nearing the performance of top-tier models like Gemini-Ultra and GPT-4. When evaluating self-consistency across 64 outputs, DeepSeekMath~7B reaches 60.9\% on the MATH benchmark. This model’s advanced mathematical reasoning stems from two pivotal innovations: firstly, an extensive use of large-scale, openly available web data, curated via a highly refined data filtering pipeline; secondly, the implementation of Group Relative Policy Optimization (GRPO)—a new variant of Proximal Policy Optimization (PPO)—which simultaneously boosts reasoning proficiency and optimizes memory consumption compared to traditional PPO approaches.
\end{abstract}

\section{Introduction}

The emergence of large language models (LLM) has fundamentally transformed the landscape of mathematical reasoning within artificial intelligence, enabling major progress on both quantitative \citep{MATH} and geometric \citep{trinh2024solving} reasoning tasks. Additionally, LLMs now play a crucial role in supporting human users faced with advanced mathematical challenges \citep{tao}. Despite these advances, leading models such as GPT-4 \citep{gpt4} and Gemini-Ultra \citep{gemini} remain closed to the public, resulting in a substantial performance gap between open-access models and these proprietary systems.

This paper presents DeepSeekMath, a language model tailored for mathematical domains that demonstrates superior mathematical performance compared to existing open-source alternatives, approaching GPT-4's accuracy on standard academic evaluations. The model’s strong results are enabled by the DeepSeekMath Corpus—a meticulously constructed pre-training dataset encompassing 120B mathematics-related tokens. The corpus is compiled from Common Crawl (CC) via a fastText-driven classification system \citep{joulin2016fasttext}. Initially, this classifier is trained on data from OpenWebMath \citep{openwebmath} labeled as positive samples, with negative samples drawn from a heterogeneous set of web documents. The trained classifier is then deployed to extract further positives from CC, which are subsequently curated and annotated by human experts. These newly labeled samples are incorporated to retrain and refine the classifier, improving its precision and recall. Evaluation reveals the substantial quality of the resulting dataset: the DeepSeekMath-Base 7B model attains a score of 64.2\% on GSM8K \citep{gsm8k} and 36.2\% on the challenging MATH benchmark \citep{MATH}, surpassing the performance of Minerva 540B \citep{minerva}. Moreover, since the DeepSeekMath Corpus is designed with multilingual content, notable gains are observed in Chinese-language mathematics assessments \citep{wei2023cmath,agieval}. We posit that our approach to compiling mathematical datasets lays valuable groundwork for future exploration, recognizing that considerable advancements remain possible.

DeepSeekMath-Base is constructed by initializing with DeepSeek-Coder-Base-v1.5 7B \citep{deepseek-coder}, as our findings suggest that commencing from a model trained on code yields superior outcomes compared to initializing with a standard large language model. Additionally, our results show that mathematical training not only augments the model’s performance on math-specific benchmarks, but also enhances results on general reasoning tasks such as MMLU \citep{mmlu} and BBH \citep{bbh}. This suggests a broader boost in reasoning capabilities beyond just mathematics.

Subsequent to pre-training, we fine-tune DeepSeekMath-Base using mathematical instruction data that incorporates chain-of-thought \citep{cot}, program-of-thought \citep{pot,pal}, and tool-augmented reasoning examples \citep{tora}. This instruction-tuned model, DeepSeekMath-Instruct 7B, surpasses all other 7B models and demonstrates performance on par with instruction-tuned open-source models at the 70B scale.

We also propose Group Relative Policy Optimization (GRPO), which modifies the standard Proximal Policy Optimization (PPO) framework \citep{schulman2017proximal}. Unlike PPO, GRPO does not rely on a critic model but instead computes the baseline through group scoring, thereby markedly reducing computational demands. Training GRPO with only a subset of English instruction fine-tuning data yields marked improvements over the already strong DeepSeekMath-Instruct. Notably, this method achieves notable gains on both in-domain benchmarks (GSM8K: 82.9\% $\rightarrow$ 88.2\%, MATH: 46.8\% $\rightarrow$ 51.7\%) and on out-of-domain mathematical tasks (e.g., CMATH: 84.6\% $\rightarrow$ 88.8\%) during RL training. Furthermore, we offer a unified analytical framework encompassing a variety of techniques including Rejection Sampling Fine-Tuning (RFT) \citep{yuan2023scaling}, Direct Preference Optimization (DPO) \citep{dpo}, PPO, and GRPO. Within this framework, these algorithms can be viewed as either direct or simplified variants of reinforcement learning. To deepen our understanding of this paradigm, we carry out a comprehensive series of experiments, exploring aspects such as online versus offline optimization, outcome versus process supervision, and single-turn versus iterative reinforcement learning. Finally, we elucidate the mechanisms underlying the RL-driven improvements in instruction-tuned model performance, and highlight avenues for further optimizing reinforcement learning within this unified framework.

\subsection{Contributions}
This work presents advances in scalable mathematical pre-training and delves into the evaluation and utilization of reinforcement learning.

\textbf{Math Pre-Training at Scale}

\begin{itemize}[topsep=0pt]
    \item We demonstrate that the open-access Common Crawl resource possesses considerable mathematical relevance. Through a carefully structured data filtering framework, we assemble the DeepSeekMath Corpus—a curated dataset comprising 120B tokens sourced from mathematically oriented web content. This dataset is nearly seven times larger than the math-specific web data employed by Minerva \citep{minerva} and exceeds the size of OpenWebMath \citep{openwebmath} by a factor of nine.

    \item Our pre-trained foundation model, DeepSeekMath-Base 7B, reaches performance levels similar to those of Minerva 540B \citep{minerva}, suggesting that sheer parameter count is not the principal determinant of mathematical reasoning strength. Robust performance can be realized by smaller models provided they are trained on rich, high-quality datasets.

    \item We provide insights drawn from our training runs in mathematics. Pre-training models on code prior to mathematical data consistently enhances their capacity to solve mathematical questions, with and without auxiliary tool support. This contributes to the discourse around the effect of code training on reasoning, and our findings indicate code pre-training does support improvements, at least in the mathematical domain.

    \item Despite the widespread practice of pre-training with arXiv material in math-focused language models, our results indicate this approach yields no substantial gains across the mathematical benchmarks considered in this study.
\end{itemize}

\textbf{Exploration and Analysis of Reinforcement Learning}

\begin{itemize}[topsep=0pt]
    \item We present Group Relative Policy Optimization (GRPO), a novel reinforcement learning approach characterized by its simplicity and efficiency. GRPO dispenses with the critic component and instead calculates the baseline using group-wise scores, thereby considerably lowering computational and memory costs in comparison to Proximal Policy Optimization (PPO).
    \item We show that, when applied solely to instruction-tuning data, GRPO considerably boosts the capabilities of our instruction-tuned model, DeepSeekMath-Instruct. In addition, we notice improved generalization on out-of-domain tasks throughout the reinforcement learning stage.
    \item We offer a cohesive framework for interpreting methods like RFT, DPO, PPO, and GRPO. Our thorough empirical analysis encompasses aspects such as online versus offline learning, outcome versus process-based supervision, and single-turn versus iterative reinforcement learning, thereby elucidating the fundamental components of this framework.
    \item Leveraging our unified framework, we investigate the underlying mechanisms contributing to the effectiveness of reinforcement learning, and identify multiple promising avenues for advancing reinforcement learning in large language models (LLMs).
\end{itemize}

\subsection{Summary of Evaluations and Metrics}

\begin{itemize}[topsep=0pt]
    \item \textbf{English and Chinese Mathematical Reasoning}: Our models are systematically evaluated on a broad spectrum of English and Chinese mathematical reasoning benchmarks, spanning tasks from elementary to collegiate level. The English benchmarks comprise GSM8K \citep{gsm8k}, MATH \citep{MATH}, SAT \citep{llemma}, OCW Courses \citep{minerva}, and MMLU-STEM \citep{mmlu}, while the Chinese suite encompasses MGSM-zh \citep{mgsm}, CMATH \citep{wei2023cmath}, Gaokao-MathCloze \citep{agieval}, and Gaokao-MathQA \citep{agieval}. Assessments include both models' capacity to autonomously compose detailed solutions (without external tools) and their proficiency at resolving questions via Python.

    On English tasks, DeepSeekMath-Base attains performance levels competitive with the proprietary Minerva 540B \citep{minerva}, outperforming all open-source base models (such as Mistral 7B \citep{mistral} and Llemma-34B \citep{llemma}) by notable margins, regardless of their pre-training regimens. Importantly, DeepSeekMath-Base also exhibits clear superiority on the Chinese evaluation suite. We attribute this, in part, to the incorporation of high-quality multilingual math data, as opposed to the English-exclusive approach adopted in prior studies \citep{minerva,llemma}. Further improvements are observed with instruction tuning and reinforcement learning: DeepSeekMath-Instruct and DeepSeekMath-RL deliver robust results, achieving above 50\% accuracy on the rigorous, competition-level MATH benchmark—a first for open-source models.
\end{itemize}

\item \textbf{Formal Mathematics}: DeepSeekMath-Base is assessed using the informal-to-formal theorem proving benchmark established in \citep{dsp_proof}, utilizing miniF2F \citep{minif2f} as the evaluation dataset and Isabelle \citep{isabelle} as the selected proof assistant. The results show that DeepSeekMath-Base achieves notable performance in few-shot autoformalization tasks.
\item \textbf{Natural Language Understanding, Reasoning, and Code}: To provide a thorough evaluation of general comprehension, reasoning, and programming skills, DeepSeekMath-Base is tested on several benchmarks: Massive Multitask Language Understanding (MMLU) \citep{mmlu}, comprising 57 multiple-choice exams from a broad range of topics; BIG-Bench Hard (BBH) \citep{bbh}, which contains 23 highly challenging tasks focused primarily on complex, multi-step reasoning; as well as HumanEval \citep{codex} and MBPP \citep{mbpp}, both standard benchmarks for assessing code generation. The process of mathematical pre-training is found to enhance both language understanding and reasoning outcomes.

\section{Math Pre-Training}

\subsection{Data Collection and Decontamination} 
This section details the methodology for assembling the DeepSeekMath Corpus from Common Crawl. Illustrated in Figure \ref{fig:data_collect}, an iterative approach is described for systematically acquiring a large mathematical dataset from Common Crawl, beginning with a seed dataset—typically a concise, high-quality set of mathematical resources. This procedure is also adaptable for constructing corpora in fields such as programming.

Initially, OpenWebMath \citep{openwebmath}, a reputable resource of mathematical texts found online, is designated as the seed dataset. Leveraging this corpus, a fastText model \citep{joulin2016fasttext} is trained to retrieve further web content with similar mathematical characteristics. For this, 500,000 entries are randomly sampled from the seed collection as positive cases, while another 500,000 web pages from Common Crawl serve as negative examples. The fastText model is trained with the help of an open-source package\footnote{\url{https://fasttext.cc}}, using the following settings: 256 for vector size, a learning rate of 0.1, up to 3-word n-grams, a minimum occurrence of 3 for each word, and 3 epochs of training. To effectively decrease the scale of Common Crawl, deduplication—both based on URLs and content similarity—is applied, narrowing down the set to 40B HTML web pages. Employing the fastText model, potentially mathematical web pages are identified within this filtered collection. These are subsequently ranked by the fastText model’s predicted relevance scores, with only the highest-ranking content retained. To determine the most appropriate data volume to preserve, pre-training is conducted on subsets with the top 40B, 80B, 120B, and 160B tokens, respectively. In the initial iteration, the process opts to select the top 40B tokens.

Following the initial phase of data collection, a significant portion of mathematical web content remains uncollected. This is primarily due to the limited diversity within the positive example set used to train the fastText model. To address this, we expand the set of mathematical web sources included in the seed corpus, aiming to enhance the model’s performance. Specifically, we first partition the entirety of the Common Crawl dataset by domain, where each domain groups all web pages sharing the same base URL. For every domain, we determine the fraction of its web pages that were gathered during the initial iteration. Domains with more than 10\% of their pages collected are designated as math-focused (for example, \url{mathoverflow.net}).

Next, we manually annotate those URLs within the selected domains that are linked to mathematical content (such as \url{mathoverflow.net/questions}). Web pages associated with these annotated URLs, which have not yet been collected, are added to our seed corpus. This method increases the number of positive samples, allowing us to retrain fastText for broader coverage in subsequent iterations. After repeating this collection process for four rounds, we assemble a dataset comprising 35.5 million mathematical web pages, containing a total of 120 billion tokens. By the fourth iteration, we observe that nearly 98\% of pages acquired in the third iteration have already been collected, prompting us to halt additional data gathering.

To prevent contamination with benchmark materials, we employ a filtering protocol aligned with \cite{deepseek-coder}, removing web pages containing content from English mathematical benchmarks such as GSM8K \citep{gsm8k} and MATH \citep{MATH}, as well as Chinese benchmarks like CMATH \citep{wei2023cmath} and AGIEval \citep{agieval}. Specifically, we discard any text block featuring a 10-gram sequence that matches verbatim any substring from the benchmark datasets. For benchmark texts shorter than 10 grams but with a minimum length of 3 grams, we apply strict exact matching to filter out compromised web pages from our math training set.

\subsection{Validating the Quality of the DeepSeekMath~Corpus}

To evaluate the relative quality of the DeepSeekMath~Corpus, we conduct pre-training experiments against several recently introduced math-focused training corpora:

\begin{itemize}[topsep=0pt]
    \item \textbf{MathPile} \citep{mathpile}: a compilation of sources totaling 8.9 billion tokens, sourced from textbooks, Wikipedia, ProofWiki, CommonCrawl, StackExchange, and predominantly (over 85\%) arXiv;
    \item \textbf{OpenWebMath} \citep{openwebmath}: a 13.6 billion-token collection of CommonCrawl web pages selected for mathematical relevance;
    \item \textbf{Proof-Pile-2} \citep{llemma}: this dataset merges OpenWebMath, AlgebraicStack (contributing 10.3B tokens of mathematical code), and arXiv publications (28.0B tokens). Consistent with \cite{llemma}, we adopt an arXiv:Web:Code composition ratio of 2:4:1 in our experiments involving Proof-Pile-2.
\end{itemize}

\subsubsection{Training Setting}
\label{sec:quality-policy}
We perform mathematical training on a general-purpose language model with 1.3B parameters, adhering to the architecture utilized by DeepSeek LLMs \citep{deepseek-llm}, and refer to this variant as DeepSeek-LLM 1.3B. Distinct models are fine-tuned individually on each respective mathematical dataset for a total of 150B tokens. The training processes make use of the efficient HAI-LLM \citep{haillm} infrastructure. Consistent with DeepSeek LLMs' training protocols, the optimization is managed using AdamW \citep{adamW}, configured with $\beta_1=0.9$, $\beta_2=0.95$, and $\mathrm{weight\_decay}=0.1$. The learning rate is governed by a multi-phase schedule: after 2,000 warmup iterations, the learning rate reaches its peak value, then declines to 31.6\% of the maximum after 80\% of total training, and subsequently tapers to 10.0\% at the 90\% mark. The peak learning rate is set at 5.3e-4, with batches consisting of 4 million tokens, and a context window of 4,000 tokens.

\subsubsection{Evaluation Results}

\textbf{The DeepSeekMath~Corpus exhibits outstanding quality, provides extensive multilingual mathematical data, and surpasses other corpora in scale.}

\begin{itemize}[topsep=0pt]
    \item \textbf{High-quality}:
    We measure downstream performance across eight mathematical benchmarks utilizing few-shot chain-of-thought prompting \cite{cot}. Table \ref{tab:corpora_comparison} demonstrates that models trained on DeepSeekMath~Corpus achieve superior results. Furthermore, Figure \ref{fig:corpora_comparisons} illustrates that the DeepSeekMath-trained model outperforms its counterpart trained on Proof-Pile-2 at the 50B-token mark (corresponding to a complete pass through Proof-Pile-2), signifying the superior average quality of DeepSeekMath~Corpus.
    \item \textbf{Multilingual}:
    DeepSeekMath~Corpus includes texts in multiple languages, with English and Chinese making up the majority of the dataset. As reported in Table \ref{tab:corpora_comparison}, models fine-tuned on DeepSeekMath~Corpus exhibit enhanced mathematical reasoning abilities in both English and Chinese. Conversely, most existing math corpora focus primarily on English, which not only limits their effectiveness but may even have negative impacts on Chinese-language mathematical reasoning tasks.
    \item \textbf{Large-scale}:
    DeepSeekMath~Corpus is substantially larger compared to other available mathematical datasets. As seen in Figure \ref{fig:corpora_comparisons}, DeepSeek-LLM 1.3B, when trained on DeepSeekMath~Corpus, demonstrates a more rapid improvement trajectory and continued gains. On the other hand, alternative baseline corpora are much more limited in size, necessitating repeated exposures during training, which leads to models quickly reaching a performance ceiling.
\end{itemize}

\subsection{Training and Evaluating DeepSeekMath-Base 7B}

This section presents DeepSeekMath-Base 7B, a foundational model designed for advanced reasoning tasks, with a particular emphasis on mathematics. The model initialization utilizes DeepSeek-Coder-Base-v1.5 7B \citep{deepseek-coder} and it undergoes training on 500B tokens. The composition of the training dataset is as follows: 56\% sourced from the DeepSeekMath Corpus, 4\% from AlgebraicStack, 10\% from arXiv, 20\% consisting of Github code, and the final 10\% comprised of natural language data from Common Crawl in both English and Chinese. The training process largely adheres to the configuration detailed in Section \ref{sec:quality-policy}, with two key exceptions: the learning rate is capped at 4.2e-4 and the batch size is set to 10 million tokens.

To thoroughly evaluate the mathematical proficiency of DeepSeekMath-Base 7B, we examine its capacity for generating independent mathematical solutions, solving tasks that require tool use, and formal theorem proving. We further supplement this analysis by assessing general capabilities such as natural language understanding, complex reasoning, and programming proficiency.

\paragraph{Mathematical Problem Solving with Step-by-Step Reasoning}

The evaluation of DeepSeekMath-Base's mathematical problem-solving is conducted using few-shot chain-of-thought prompting \citep{cot} across eight distinct benchmarks in both English and Chinese. These datasets include tests of quantitative reasoning—such as GSM8K \citep{gsm8k}, MATH \citep{MATH}, and CMATH \citep{wei2023cmath}—as well as multiple-choice formats like MMLU-STEM \citep{mmlu} and Gaokao-MathQA \citep{agieval}. Collectively, these benchmarks represent a broad spectrum of mathematical domains and difficulty levels, ranging from elementary concepts to those encountered in university curricula.

As detailed in Table \ref{tab:base_math_cot}, among publicly available base models, DeepSeekMath-Base 7B achieves the highest scores across all eight evaluation benchmarks, outperforming established models such as the widely adopted Mistral 7B \citep{mistral} and the newly introduced Llemma 34B \citep{llemma}, the latter having been additionally trained on Proof-Pile-2 \citep{llemma} for mathematical proficiency. In particular, DeepSeekMath-Base attains over a 10\% absolute improvement compared to previous open-source base models on the MATH competition-level dataset, and even exceeds the performance of Minerva 540B \citep{minerva}—a proprietary model with 77 times the parameter count, developed from PaLM \citep{palm} with further domain-specific training.

\paragraph{Mathematical Problem Solving with Tool Use} To assess capabilities in program-assisted mathematical reasoning, we utilize few-shot program-of-thought prompting \citep{pot,pal} on the GSM8K and MATH benchmarks. Here, models approach each problem by scripting a Python solution, optionally leveraging packages like \textit{math} and \textit{sympy} for advanced calculations. The accuracy of the answer is determined by the output from executing these scripts. According to Table \ref{tab:base_math_pot_proof}, DeepSeekMath-Base 7B achieves results superior to the previous best-performing Llemma 34B.

\paragraph{Formal Mathematics}

Automating formal proofs has become increasingly relevant for ensuring the correctness and dependability of mathematical arguments and streamlining their construction. For evaluation, we apply DeepSeekMath-Base 7B to the informal-to-formal proving task described in \citep{dsp_proof}, which involves creating a formal proof from an informal problem description, a formalized version of the statement, and an informal argument. We conduct this evaluation on the miniF2F \citep{minif2f} benchmark—targeting Olympiad-level mathematics—generating Isabelle formal proofs with few-shot prompting. Consistent with the protocol in \cite{dsp_proof}, the model is used to propose proof outlines, which are then supplemented by Sledgehammer \citep{sledgehammer}, an automated prover that completes the remaining steps. Table \ref{tab:base_math_pot_proof} indicates that DeepSeekMath-Base 7B demonstrates notable effectiveness in automated formalization.

\paragraph{Natural Language Understanding, Reasoning, and Code}

We benchmark the model’s natural language understanding on MMLU \citep{mmlu}, reasoning abilities on BBH \citep{bbh}, and coding proficiency using HumanEval \citep{codex} and MBPP \citep{mbpp}. As reported in Table \ref{tab:understanding_reasoning_code}, DeepSeekMath-Base 7B achieves marked improvements on both MMLU and BBH over its predecessor, DeepSeek-Coder-Base-v1.5 \citep{deepseek-coder}, indicating that math-specific training bolsters both language comprehension and logical reasoning. Furthermore, through the integration of code tokens during continued training, DeepSeekMath-Base 7B sustains the performance level of DeepSeek-Coder-Base-v1.5 across the coding tasks. Collectively, DeepSeekMath-Base 7B consistently surpasses Mistral 7B \citep{mistral} on all three benchmarks involving reasoning and programming skills.

\section{Supervised Fine-Tuning}

\subsection{SFT Data Curation}

We compile a comprehensive instruction-tuning dataset for mathematics, encompassing English and Chinese tasks across various mathematical domains and difficulty levels. Each problem is coupled with a corresponding solution provided in the format of chain-of-thought (CoT) \citep{cot}, program-of-thought (PoT) \citep{pot,pal}, or tool-integrated reasoning \citep{tora}. In total, the training set contains 776K examples.

\begin{itemize}[topsep=0pt]
    \item \textbf{English mathematical datasets}: We enrich GSM8K and MATH with tool-integrated solution annotations, and incorporate a subset of MathInstruct \citep{MathInstruct} alongside the Lila-OOD training split \citep{lila}, in which problems are resolved using either CoT or PoT approaches. The English dataset spans a wide spectrum of mathematical topics such as algebra, probability, number theory, calculus, and geometry.
    \item \textbf{Chinese mathematical datasets}: We gather a collection of Chinese K-12 math exercises representing 76 sub-fields (e.g., linear equations), with answers detailed in both CoT and tool-integrated reasoning structures.
\end{itemize}

\subsection{Training and Evaluating DeepSeekMath-Instruct 7B}

This section details the training regimen for DeepSeekMath-Instruct 7B, which is refined through instruction tuning built on DeepSeekMath-Base. Training samples are concatenated at random up to a context window of 4K tokens. The optimization employs 500 training steps, with a batch size of 256 and a fixed learning rate set at 5e-5.

Model assessment encompasses mathematical capabilities with and without access to external tools, employing four quantitative reasoning benchmarks in both English and Chinese. We compare the results of our approach with other state-of-the-art models contemporaneous to our work:

\begin{itemize}[topsep=0pt]
    \item \textbf{Closed-source models} considered are: (1) the GPT model series, particularly GPT-4 \citep{gpt4} and GPT-4 Code Interpreter~\footnote{\url{https://openai.com/blog/chatgpt-plugins\#\#code-interpreter}}, (2) Gemini Ultra and Pro \citep{gemini}, (3) Inflection-2 \citep{inflection-2}, (4) Grok-1~\footnote{\url{https://x.ai/model-card}}, as well as recently introduced Chinese models, including (5) Baichuan-3~\footnote{\url{https://www.baichuan-ai.com}}, and (6) the newest addition to the GLM lineup, GLM-4~\footnote{\url{https://open.bigmodel.cn/dev/api\#glm-4}} from the GLM family \citep{glm}.
\end{itemize}

These models are generally designed for broad applications and typically undergo various alignment processes.
\item \textbf{Open-source models} encompass both general-purpose and math-specialized variants, such as (1) DeepSeek-LLM-Chat 67B \citep{deepseek-llm}, (2) Qwen 72B \citep{qwen}, (3) SeaLLM-v2 7B \citep{seallm}, and (4) ChatGLM3 6B \citep{chatglm3}. The category further includes models optimized for mathematical tasks: (5) InternLM2-Math 20B~\footnote{\url{https://github.com/InternLM/InternLM-Math}}, an adaptation of InternLM2 subjected to specialized math pre-training and subsequent instruction tuning; (6) Math-Shepherd-Mistral 7B, which integrates PPO training \citep{schulman2017proximal} into Mistral 7B \citep{mistral} guided by a process-supervised reward system; (7) the WizardMath collection \citep{wizardmath}, enhancing mathematical capabilities of Mistral 7B and Llama-2 70B \citep{llama2} through evolve-instruct—a form of instruction tuning utilizing instructions generated by AI—and PPO optimization, with data predominantly from GSM8K and MATH; (8) MetaMath 70B \citep{metamath}, comprising Llama-2 70B fine-tuned using enriched GSM8K and MATH corpora; (9) ToRA 34B \cite{tora}, which leverages CodeLlama 34B and is tailored for mathematical reasoning involving tool integration; and (10) MAmmoTH 70B \citep{MathInstruct}, where Llama-2 70B is refined via instruction tuning with MathInstruct.

Table \ref{tab:sft_rl_math} illustrates that in the tool-free evaluation scenario, DeepSeekMath-Instruct 7B delivers robust performance in detailed, stepwise reasoning. Most prominently, on the competition-grade MATH dataset, it outperforms every open-source alternative as well as the vast majority of proprietary models (such as Inflection-2 and Gemini Pro) by a margin of at least 9\% in absolute accuracy, including those that are significantly larger (e.g., Qwen 72B) or have received targeted enhancements through reinforcement learning techniques for math (e.g., WizardMath-v1.1 7B). Although DeepSeekMath-Instruct is comparable to leading Chinese proprietary systems like GLM-4 and Baichuan-3 on MATH, it still lags behind GPT-4 and Gemini Ultra.

When evaluating in the setting where systems can jointly leverage natural language reasoning and programmatic tools, DeepSeekMath-Instruct 7B achieves nearly 60\% accuracy on MATH, outstripping all publicly available open-source models. On other assessment datasets, our system remains competitive with DeepSeek-LLM-Chat 67B, the previous top performer, despite being one-tenth the size.
\section{Reinforcement Learning}

\subsection{Group Relative Policy Optimization} The reinforcement learning (RL) paradigm has demonstrated substantial benefit for further enhancing LLMs’ capacity for mathematical problem solving following Supervised Fine-Tuning (SFT) \citep{wang2023math,wizardmath}. In this section, we present Group Relative Policy Optimization (GRPO), an RL method that is both efficient and effective.

\subsubsection{From PPO to GRPO} Proximal Policy Optimization (PPO) \citep{schulman2017proximal} stands as a prominent actor-critic method for RL, widely employed in the fine-tuning phase of LLMs via RL \citep{ouyang2022training}. PPO improves LLMs by maximizing a surrogate objective given by:

\begin{equation}
\footnotesize
    \mathcal{J}_{PPO}(\theta) = \mathbb{E}{[q \sim P(Q), o \sim \pi_{\theta_{old}}(O|q)]} \frac{1}{|o|} \sum_{t=1}^{|o|} \min \left[ \frac{\pi_\theta(o_{t} | q, o_{<t})}{\pi_{\theta_{old}}(o_{t} | q, o_{<t})} A_{t}, \text{clip} \left( \frac{\pi_\theta(o_{t} | q, o_{<t})}{\pi_{\theta_{old}}(o_{t} | q, o_{<t})}, 1 - \varepsilon, 1 +

Here, $\pi_{\theta}$ and $\pi_{\theta_{old}}$ denote the current and previous policy models, respectively, and $q$ and $o$ represent questions and outputs that are drawn from the question dataset as well as from the previous policy $\pi_{\theta_{old}}$. The parameter $\varepsilon$ is a clipping hyper-parameter used in PPO to help stabilize the optimization process. The term $A_t$ corresponds to the advantage, which is obtained through Generalized Advantage Estimation (GAE) \citep{gae}, leveraging the sequence of rewards $\{r_{\ge t}\}$ and a parameterized value function $V_{\psi}$. Consequently, PPO requires training a value function in addition to the policy model, and to curb reward model over-optimization, the standard methodology is to incorporate a per-token KL penalty relative to a reference model within the reward at each token position \citep{ouyang2022training}, formally expressed as

\begin{equation}
    r_{t} = r_\varphi(q, o_{\le t}) - \beta \log\frac{\pi_{\theta}(o_{t}|q, o_{<t})}{\pi_{ref}(o_{t}|q, o_{<t})},
\label{eq:PPO-reward}
\end{equation}

where $r_\varphi$ denotes the reward model, $\pi_{ref}$ refers to the reference model (commonly the initial SFT model), and $\beta$ indicates the weight of the KL penalty.

As the value function in PPO is typically implemented as a model with similar capacity to the policy model, it introduces considerable computational and memory demands. Additionally, in RL-based training, this value function serves as a baseline in advantage estimation for variance reduction. In the context of LLMs, however, reward scores are usually assigned only to the final token of the generated sequence, making accurate value function training per token challenging. To overcome these issues, we introduce Group Relative Policy Optimization (GRPO), as illustrated in Figure \ref{fig:grpo}, which eliminates the need for training a separate value function as in PPO. Instead, GRPO computes a baseline using the mean reward across multiple sampled responses to the same question. Specifically, for any given question $q$, a collection of outputs $\{o_1, o_2, \cdots, o_G\}$ is generated via the old policy $\pi_{\theta_{old}}$, and policy optimization is performed by maximizing the following objective:

\begin{equation}
\footnotesize
\begin{split}
    \mathcal{J}_{GRPO}(\theta) &= \mathbb{E}{[q \sim P(Q), \{o_i\}_{i=1}^G \sim \pi_{\theta_{old}}(O|q)]}  \\
    & \frac{1}{G}\sum_{i=1}^G\frac{1}{|o_i|} \sum_{t=1}^{|o_i|} \left\{ \min \left[ \frac{\pi_\theta(o_{i,t} | q, o_{i,<t})}{\pi_{\theta_{old}}(o_{i,t} | q, o_{i,<t})} \hat{A}_{i,t}, \text{clip} \left( \frac{\pi_\theta(o_{i,t} | q, o_{i,<t})}{\pi_{\theta_{old}}(o_{i,t} | q, o_{i,<t})}, 1 - \varepsilon, 1 + \varepsilon \right)  \hat{A}_{i,t} \right] - \beta \mathbb{D}_{KL}\left[\pi_{\theta} || \pi_{ref}\right]\right\} ,
\end{split}
\label{eq:GRPO-obj}
\end{equation}

where $\varepsilon$ and $\beta$ are hyper-parameters, and $\hat{A}_{i,t}$ stands for the advantage, which is computed solely based on the relative rewards among group members (elaborated in subsequent subsections). This group-relative computation of advantages in GRPO fits naturally with the comparative training schemes of reward models, since these models are generally optimized using datasets that compare alternative responses to identical questions. Furthermore, GRPO diverges from the practice of including a KL penalty within the reward; instead, it applies a KL divergence regularization term directly to the loss function, thereby simplifying the computation of $\hat{A}_{i,t}$. Unlike the KL penalty used in (\ref{eq:PPO-reward}), the KL divergence in GRPO is estimated via the following unbiased estimator \citep{kl_approx}:

\begin{equation}
\

\begin{algorithm}[t]
  \small
  \caption{Iterative Group Relative Policy Optimization}
  \textbf{Input} initial policy model $\pi_{\theta_{\text{init}}}$; reward models $r_\varphi$; task prompts $\mathcal{D}$; 
  hyperparameters $\varepsilon$, $\beta$, $\mu$
  \begin{algorithmic}[1]
    \State Initialize policy model: $\pi_\theta \leftarrow \pi_{\theta_{\text{init}}}$
    \For{iteration = 1, \dots, I}
       \State Set reference model $\pi_{ref} \leftarrow \pi_{\theta}$
      \For{step = 1, \dots, M}
      \State Draw a batch $\mathcal{D}_b$ from $\mathcal{D}$
      \State Assign current policy to previous: $\pi_{\theta_{old}} \leftarrow \pi_{\theta}$ 
      \State For each query $q \in \mathcal{D}_b$, generate $G$ responses $\{o_i\}_{i=1}^G \sim \pi_{\theta_{old}} (\cdot \mid q) $
      \State For each generated $o_i$, evaluate rewards $\{r_i\}_{i=1}^{G}$ using $r_{\varphi}$
      \State Calculate group-relative advantages $\hat{A}_{i,t}$ at token $t$ for every $o_i$
      \For{GRPO iteration = 1, \dots, $\mu$}
        \State Update $\pi_{\theta}$ by maximizing the Group-Relative Policy Optimization target (Equation \ref{eq:GC-GRPO})
      \EndFor
    \EndFor \State Continuously train $r_\varphi$ via replay mechanism. 
    \EndFor 
  \end{algorithmic}
  \textbf{Output} $\pi_\theta$
  \label{alg:iter-grpo}
\end{algorithm}

\subsubsection{Outcome Supervision RL with GRPO} In this framework, for each question $q$, a collection of candidate outputs $\{o_1, o_2, \cdots, o_G\}$ is produced using the previous policy $\pi_{\theta_{old}}$. Each candidate output is assigned a reward by the reward model, yielding a reward vector $\mathbf{r} = \{r_1, r_2, \cdots, r_G\}$. These raw rewards are standardized by subtracting the mean of $\mathbf{r}$ and dividing by its standard deviation. Outcome-based supervision assigns the resulting normalized value to each token in output $o_i$ as its advantage, so that $\hat{A}_{i, t} = \widetilde{r}_i = \frac{r_i- {\rm mean}(\mathbf{r})}{{\rm std}(\mathbf{r})}$ for all $t$. Policy parameters are updated by maximizing the target specified in equation (\ref{eq:GRPO-obj}).

\subsubsection{Process Supervision RL with GRPO} In contrast to outcome supervision, which yields a single reward after each output, process supervision provides intermediate rewards at every step, which is especially valuable for guiding policies in multifaceted mathematical reasoning, as described in \cite{wang2023math}. Specifically, for a question $q$ and a set of $G$ outputs $\{o_1, o_2, \cdots, o_G\}$, a process-based reward model produces rewards for every reasoning step: $\mathbf{R} = \{ \{r_1^{index(1)},\cdots,r_1^{index(K_1)}\}, \cdots,  \{r_G^{index(1)},\cdots,r_G^{index(K_G)}\} \}$, with $index(j)$ marking the end position of the $j$-th step, and $K_i$ the number of steps in $o_i$. Each step reward is normalized by the global mean and standard deviation: $\widetilde{r}_i^{index(j)} = \frac{r_i^{index(j)} - {\rm mean(\mathbf{R})}}{{\rm std(\mathbf{R})}}$. 
Process supervision then determines the advantage at token $t$ as the sum of subsequent normalized step rewards: $\hat{A}_{i, t} = \sum_{index(j) \ge t} \widetilde{r}_i^{index(j)}$. Policy optimization proceeds by maximizing the same objective from equation (\ref{eq:GRPO-obj}).

\subsubsection{Iterative RL with GRPO} During the reinforcement learning procedure, the reward model previously used may no longer provide adequate supervision for the evolving policy model. To address this, we investigate an iterative approach to RL employing GRPO. As described in Algorithm \ref{alg:iter-grpo}, the iterative GRPO framework periodically creates new datasets for the reward model utilizing samples generated by the latest policy model. The reward model is updated incrementally using a replay buffer that includes 10\% of data from prior iterations. Subsequently, we update the reference model to the most recent policy model, after which the policy model is further refined under the guidance of the updated reward model.

\subsection{Training and Evaluating DeepSeekMath-RL}

We employ DeepSeekMath-Instruct 7B as the foundation for our reinforcement learning experiments. RL training leverages chain-of-thought-style items drawn from GSM8K and MATH within the SFT dataset, comprising roughly 144K problems. Other SFT question types are intentionally omitted in order to evaluate RL's effect on test sets that are not represented during RL training. The reward model’s training data is assembled as in \citep{wang2023math}. The initial reward model is trained on DeepSeekMath-Base 7B, using a learning rate of 2e-5. The policy model in GRPO utilizes a learning rate of 1e-6, and the KL divergence coefficient is set to 0.04. We sample $64$ completions per question. Sequences are capped at a length of 1024 tokens, and the batch size during training is 1024. For each round of exploration, the policy model undergoes only a single parameter update. DeepSeekMath-RL 7B is benchmarked using the same protocols as DeepSeekMath-Instruct 7B. Within this evaluation, GSM8K and MATH chain-of-thought tasks are categorized as in-domain, while all remaining benchmarks are treated as out-of-domain assessments.

Table \ref{tab:sft_rl_math} presents the outcomes of both open-source and closed-source models, evaluating their performance in chain-of-thought and tool-augmented reasoning tasks on English and Chinese datasets. Key observations include: 1) DeepSeekMath-RL 7B achieves accuracy rates of 88.2\% on GSM8K and 51.7\% on MATH benchmarks when employing chain-of-thought reasoning. These results not only surpass all open-source models ranging from 7B to 70B in size but also outperform the majority of closed-source models. 2) Importantly, DeepSeekMath-RL 7B is fine-tuned exclusively on chain-of-thought-structured instruction data from GSM8K and MATH, using DeepSeekMath-Instruct 7B as its foundation. Even though its training data is restricted to these formats, DeepSeekMath-RL 7B consistently exceeds the performance of DeepSeekMath-Instruct 7B across all metrics, underscoring the advantage conferred by reinforcement learning.

\section{Discussion}
This section discusses insights derived from our pre-training and reinforcement learning (RL) experiments.

\subsection{Lessons Learnt in Pre-Training}

Here, we elaborate on insights gained during the pre-training phase. Unless indicated otherwise, we follow the experimental setup provided in Section \ref{sec:quality-policy}. For this analysis, the DeepSeekMath Corpus refers to a collection comprising 89B tokens, corresponding to the second round of data assembly.

\subsubsection{Code Training Benefits Mathematical Reasoning}

While it is widely speculated—though seldom empirically validated—that code training boosts reasoning abilities, our work attempts to address this notion within the context of mathematical reasoning: code training enhances a model’s mathematical reasoning capabilities, both with and without the assistance of external tools.

To evaluate the influence of code training on mathematical reasoning, we explored the following setups: two-stage training and one-stage training.

\textbf{Two-Stage Training}

\begin{itemize}[topsep=0pt]
    \item \textbf{Code Training for 400B Tokens $\rightarrow$ Math Training for 150B Tokens}: DeepSeek-LLM 1.3B is initially trained on 400B code tokens, then further trained on 150B tokens of mathematical data;
    \item \textbf{General Training for 400B Tokens $\rightarrow$ Math Training for 150B Tokens}: For comparison, we replace code tokens with general-domain tokens (sampled from a comprehensive corpus compiled by DeepSeek-AI) in the initial pretraining phase to analyze whether code-based pretraining provides any special benefit for enhancing mathematical reasoning compared to generic textual data.
\end{itemize}

\textbf{One-Stage Training}

\begin{itemize}[topsep=0pt]
    \item \textbf{Math Training for 150B Tokens}: The model undergoes training exclusively on 150B tokens from mathematics datasets;
    \item \textbf{Training on a mixture of 400B Code Tokens and 150B Math Tokens}: Since subsequent math training after code pretraining impairs coding abilities, we explore a joint one-stage training regime, in which code and math tokens are presented together. This setting evaluates whether such a blend both enhances mathematical reasoning and minimizes catastrophic forgetting.
\end{itemize}

\paragraph{Results}
As illustrated by Table \ref{tab:code-math} and Table \ref{tab:code-math-general-eval}, different training regimens produce varied downstream outcomes.

Code-focused pretraining significantly aids the model's capacity for program-aided mathematical problem solving in both sequential (two-stage) and integrated (one-stage) scenarios. Results in Table \ref{tab:code-math} indicate that pretraining solely on code already greatly strengthens the model's performance on GSM8K and MATH benchmarks when Python tools are allowed. An additional math-specialized training phase further amplifies these improvements. Furthermore, the one-stage blended training approach, which uses both code and math tokens, effectively counters the catastrophic forgetting issue seen in sequential setups. This method enhances both programming ability (see Table \ref{tab:code-math-general-eval}) and program-supported mathematical reasoning (see Table \ref{tab:code-math}).

Beyond tool-assisted reasoning, code pretraining also yields noticeable gains. When code tokens are used in the first phase of two-stage training, even before fine-tuning with math data, the model already shows improvement. This sequence of training leads to higher overall mathematical performance following the second phase. By contrast, combining code and math tokens in a single training stage reduces the model’s effectiveness on pure mathematical reasoning tasks without tool involvement. One possible explanation is that the relatively modest model size of DeepSeek-LLM 1.3B constrains its ability to fully learn from both code and math data at once.

\subsubsection{ArXiv Papers Exhibit Limited Efficacy in Enhancing Mathematical Reasoning}

ArXiv papers frequently constitute a significant portion of mathematical pre-training corpora \citep{minerva,gpt-f,llemma,mathpile}, yet there has been little rigorous investigation into their precise effect on mathematical reasoning skills. Surprisingly, our empirical results indicate that incorporating arXiv papers does not contribute meaningfully to improvements in mathematical reasoning. We conduct experiments using language models of varying scales, specifically DeepSeek-LLM 1.3B and DeepSeek-Coder-Base-v1.5 7B \citep{deepseek-coder}, and utilize arXiv datasets subjected to different preprocessing workflows:

\begin{itemize}[topsep=0pt]
    \item \textbf{MathPile} \citep{mathpile}: a corpus of 8.9B tokens curated through extensive cleaning and filtering procedures, with more than 85\% derived from scientific arXiv articles;
    \item \textbf{ArXiv-RedPajama} \citep{redpajama}: a 28.0B-token dataset comprised of the complete collection of arXiv LaTeX sources, stripped of preambles, comments, macros, and bibliographic material.
\end{itemize}

We independently train DeepSeek-LLM 1.3B over 150B tokens and DeepSeek-Coder-Base-v1.5 7B over 40B tokens on each arXiv-based dataset. The results demonstrate that exclusive training on arXiv corpora fails to yield any significant advancement—and in some cases results in a decline—in mathematical reasoning ability, as measured by the suite of mathematical benchmarks in our analysis. These include quantitative reasoning tasks such as GSM8K and MATH (see Table \ref{tab:arxiv-cot}), multiple-choice assessments like MMLU-STEM (see Table \ref{tab:arxiv-cot}), as well as formal mathematics evaluations like miniF2F (see Table \ref{tab:arxiv_atp}).

Nonetheless, it is important to note that these findings are subject to several limitations and do not represent a definitive verdict. In particular, our investigation has not addressed the following aspects:

\begin{itemize}[topsep=0pt]
    \item The role of arXiv-derived data in specialized mathematical tasks absent from our study, such as the informalization process, where formal statements or proofs are rendered into informal language;
    \item The interaction effects when arXiv data is mixed with other data sources;
    \item Whether the advantages of training on arXiv papers might emerge only with much larger-scale models.
\end{itemize}

Therefore, more in-depth research is needed to fully understand the utility of arXiv papers in this context, which we leave as a direction for future work.

\subsection{Insights of Reinforcement Learning}
\subsubsection{Towards to a Unified Paradigm} This section seeks to present a unified analytical framework for various training techniques—including SFT, RFT, DPO, PPO, and GRPO—and to experiment with factors influencing the proposed paradigm. Typically, for these methods, the gradient with respect to the parameter $\theta$ can be formalized as:

\begin{equation}
    \nabla_{\theta}\mathcal{J}_{\textcolor{red}{\mathcal{A}}}(\theta) = \mathbb{E}[\underbrace{(q,o) \sim \textcolor{red}{\mathcal{D}}}_{Data \ Source}]\left( \frac{1}{|o|} \sum_{t=1}^{|o|}  \underbrace{GC_{{\mathcal{A}}}(q, o, t, \textcolor{red}{\pi_{{rf}}})}_{Gradient \ Coefficient}  \nabla_{\theta}\log \pi_{\theta}(o_t | q, o_{<t})\right).
\label{eq:objective}
\end{equation}

The above formulation is structured around three principal components: 1) \textit{Data Source $\mathcal{D}$}, which specifies the set of training samples; 2) \textit{Reward Function $\pi_{{rf}}$}, serving as the origin of supervision or reward; and 3) \textit{Algorithm $\mathcal{A}$}, which integrates data and reward signals to compute the gradient coefficient $GC$, quantifying how strongly each instance is reinforced or penalized. We interpret and contrast prominent training techniques within this framework:

\begin{itemize}
    \item  \textbf{Supervised Fine-tuning (SFT)}: SFT adjusts pretrained models by additional training on datasets curated by humans.
    \item \textbf{Rejection Sampling Fine-tuning (RFT)}: RFT enhances models trained with SFT by continuing training on a subset of generated responses, filtered based on answer correctness, that were sampled from the SFT model using the same set of queries.
    \item \textbf{Direct Preference Optimization (DPO)}: DPO builds upon the SFT baseline, refining model behavior through additional fine-tuning with synthetically paired outputs and a pairwise loss specific to DPO.
    \item \textbf{Online Rejection Sampling Fine-tuning (Online RFT)}: Unlike conventional RFT, Online RFT starts from the SFT initialization but collects augmented outputs on-the-fly from the live policy for continued fine-tuning.
    \item \textbf{PPO/GRPO}: These methods also initialize policies using SFT and update the model by reinforcing outputs drawn from the current policy during online learning.
\end{itemize}

The elements constituting these methods are outlined in Table \ref{tab:data_policy}. For an extensive derivation, see Appendix \ref{app:analysis-rl}.

\paragraph{Observation about Data Source} We categorize data sources into online and offline sampling. Online sampling implies that the data used for training originates from real-time exploration using the current training policy, whereas offline sampling indicates the data is derived from outputs sampled by the initial SFT model. RFT and DPO utilize an offline approach, while Online RFT and GRPO operate with online sampling.

Referring to Figure \ref{fig:rl-analysis}, our experiments reveal that Online RFT consistently surpasses RFT on both evaluation benchmarks. During the early phases of training, Online RFT achieves results comparable to RFT, but as training proceeds, it demonstrates a clear and increasing advantage, highlighting the efficacy of online approaches. This observation aligns with intuition: at the outset, the policy actor remains similar to the SFT model, so the data samples reflect only minor discrepancies. As training advances, however, the policy diverges from the initial SFT, causing online-sampled data to exhibit more substantial differences—at which point real-time sampling becomes particularly beneficial.

\paragraph{Observation about Gradient Coefficient} The method by which each algorithm translates the reward signal into a gradient coefficient is central to updating model parameters. In our setup, the reward function is either `Rule'—which assesses response quality via answer correctness—or `Model', where a trained reward model assigns a score to each response, with this model's data being labeled according to the rule-based criterion. Equations \ref{eq:GC-RFT} and \ref{eq:GC-GRPO} underscore a fundamental distinction between GRPO and Online RFT: GRPO dynamically modulates its gradient coefficient according to the reward produced by the reward model, thus enabling varied degrees of positive and negative reinforcement. Conversely, Online RFT does not employ this mechanism; it refrains from penalizing incorrect outputs and applies the same reinforcement magnitude to all correct responses.

As illustrated in Figure \ref{fig:rl-analysis}, GRPO consistently outperforms online RFT, underscoring the effectiveness of modulating the coefficients for positive and negative gradients. Moreover, GRPO+PS yields better outcomes than GRPO+OS, which demonstrates the advantage of implementing finely-tuned, step-dependent gradient adjustments. We further examine the effect of iterative reinforcement learning by applying two iterations during our experiments. Figure \ref{fig:iter-rl} reveals that this iterative RL approach considerably enhances performance, with the most notable improvement occurring after the initial iteration.

\subsubsection{Why RL Works?} This work leverages reinforcement learning on a selected subset of instruction tuning data, resulting in notable performance gains over the instruction tuning baseline. To elucidate the mechanisms behind this improvement, we analyze the Pass@K and Maj@K accuracies for both the Instruct and RL-trained models on two benchmark datasets. As depicted in Figure \ref{fig:maj-pass}, reinforcement learning increases Maj@K, but does not affect Pass@K. This suggests that the main benefit of RL lies in increasing the reliability of the output distribution—specifically, \textbf{the observed gains arise from elevating correct responses into the TopK set, as opposed to boosting the model’s intrinsic capability.} Likewise, prior work \citep{wang2023making} documents a \textbf{misalignment problem} with reasoning tasks in SFT models, and suggests that performance can be enhanced through preference alignment strategies \citep{yuan2023rrhf,song2023preference,wang2023making}.

\subsubsection{How to Achieve More Effective RL?}
\label{sec:effec_rl}
Our findings indicate that RL is highly effective for mathematical reasoning scenarios. In addition, we introduce a unified framework to conceptualize various representative training strategies, classifying all such approaches as either forms of direct RL or as simplified variants. Within this unified paradigm, three primary elements are identified—Data Source, Algorithm, and Reward Function—as outlined in Equation \ref{eq:objective}. We propose several promising future research avenues regarding these three facets.

\paragraph{Data Source} The data source forms the foundational input for any training method. In the RL setting discussed here, this specifically refers to unlabeled question instances paired with responses generated by the policy model. This study restricts its data to instruction tuning prompts and employs straightforward nucleus sampling for output generation. This limitation may partly explain why improvements from our RL methodology are confined to the Maj@K metric. Future work will investigate the applicability of our RL framework to out-of-distribution questions and explore the integration of \textbf{advanced sampling (decoding) strategies}, such as tree-search-based approaches \citep{yao2023tree}. Furthermore, leveraging \textbf{efficient inference techniques} \citep{xia-etal-2023-speculative,leviathan2023fast,kwon2023efficient,xia2024unlocking}—which impact the effectiveness of policy exploration—will be crucial for subsequent improvements.

\paragraph{Algorithms} Algorithms utilize both data and reward signals to determine the gradient coefficient for updating model parameters. According to Equation \ref{eq:objective}, contemporary approaches largely \textbf{TRUST} the guidance of the reward function to modulate the conditional probability of specific tokens. Nonetheless, the accuracy of the reward signal cannot be guaranteed in all scenarios, particularly in tasks of high complexity. As an illustration, even meticulously curated datasets such as PRM800K \citep{lightman2023let}, which benefit from expert annotation, still exhibit a notable error rate—around 20\% of the annotations are incorrect\footnote{\url{https://github.com/openai/prm800k/issues/12\#issuecomment-1728491852}}. Consequently, it becomes essential to investigate reinforcement learning strategies designed to maintain robustness in the presence of noisy reward signals. We posit that \textbf{WEAK-TO-STRONG} alignment methods \citep{burns2023weak} have the potential to fundamentally reshape the architecture of learning algorithms.

\paragraph{Reward Function} The reward function delivers the core signal that drives training. In reinforcement learning, this is typically instantiated as a neural reward model. We identify three crucial avenues for future advancement in reward model design: 1) \textbf{Enhancing the generalization capacity of reward models.} Reward models should generalize adeptly to novel, out-of-distribution queries and sophisticated model outputs; otherwise, reinforcement learning may serve only to stabilize language model output distributions without imparting substantial capability gains; 2) \textbf{Incorporating reward model uncertainty.} The explicit quantification of uncertainty could provide a critical interface connecting less reliable reward models with weak-to-strong learning paradigms; 3) \textbf{Developing efficient, high-quality process reward models} that enable the provision of granular training signals throughout multi-step reasoning tasks \citep{lightman2023let,wang2023math}.

\section{Conclusion, Limitation, and Future Work}
This work introduces DeepSeekMath, a model that surpasses all publicly available open-source counterparts on the competition-level MATH benchmark, coming close to the proficiency demonstrated by proprietary systems. DeepSeekMath~builds upon DeepSeek-Coder-v1.5 7B as its foundation and is continually trained over 500B tokens, including a substantial fraction—120B math-related tokens—extracted from Common Crawl. Our thorough ablation analysis highlights that content from web pages serves as a rich source of high-quality mathematical data, while the expected advantages from arXiv are limited. We present Group Relative Policy Optimization (GRPO), an extension of Proximal Policy Optimization (PPO), capable of substantially enhancing mathematical reasoning with reduced memory demands. Empirical evidence indicates that GRPO remains beneficial even as DeepSeekMath-Instruct 7B achieves high marks on established benchmarks. Additionally, we propose a cohesive framework to contextualize various methodologies and delineate promising avenues to advance reinforcement learning efficiency.

Despite its strong quantitative reasoning abilities, DeepSeekMath~demonstrates comparatively weaker performance in tasks involving geometry and formal proof, especially in comparison to closed-source models. As observed in preliminary assessments, the model struggles with triangle and ellipse problems, potentially signaling a bias introduced during data selection in both the pre-training and fine-tuning phases. Moreover, model size constraints mean DeepSeekMath~lags behind GPT-4 in few-shot learning scenarios; unlike GPT-4, which benefits from few-shot prompts, DeepSeekMath~displays nearly equivalent results under zero-shot and few-shot settings. Future work will focus on refining our data curation strategies to develop even more robust pre-training datasets. Furthermore, we will pursue the research directions outlined in Section~\ref{sec:effec_rl} to enhance the reinforcement learning capabilities of large language models.

\end{document}